\subsection{Exact identities and the metric mass flux}
\label{sec:ah-algebra}

We give the differential identities used in the deformation, including
their dependence on the spatially varying lapse profile.  The argument
in this subsection is local until the asymptotic assumptions in
Proposition~\ref{prop:ah-alg-mass-passage} are imposed.  In particular,
none of the identities requires a solution of the deformation equations
to have been constructed in advance.

\paragraph{Notation.}
All gradients, contractions, and divergences in the differential
identities below refer to \(g\).  We identify vectors and covectors using
\(g\) when they occur in the same formula.  Let \(t,f\) be smooth,
let \(l(t,x)>0\), and write
\[
 p=\partial_t\log l,\qquad
 A=\mathrm d_x\log l,\qquad
 B=\mathrm d\log l=A+p\,\mathrm dt.
\]
Here \(\mathrm d_x\) holds \(t\) fixed, whereas \(\mathrm d\) differentiates
the composition \(x\mapsto l(t(x),x)\). Use \(l=\lambda l_0\)
from Section~\ref{subsec:ah-profiles} and \(L=e^{2t}l_0\)
from Section~\ref{subsec:ah-equations}. Set
\begin{gather*}
 D=1+l^2|\mathrm df|^2,\qquad d=D^{-1},\qquad
 w=\frac{l\nabla f}{\sqrt D},\qquad v=|w|^2=1-d,\qquad a=\sqrt v,\\
 H^f=\frac{l\nabla^2 f}{\sqrt D},\qquad
 \overline g=g+l^2\mathrm df\otimes\mathrm df,\qquad
 \widehat g=e^{4t}\overline g,\\
 U=l\sqrt D,\qquad u=e^{2t}U=\lambda L\sqrt D,\qquad
 Z=t+\frac18\log D .
\end{gather*}
At a point where \(a>0\), put \(e=w/a\), and define
\[
 A_j=I+(j-1)e\otimes e,\qquad
 \chi=\frac{4d}{4+pv},\qquad
 S=K+H^f+w\otimes A+A\otimes w,\qquad
 F=\operatorname{tr}_{A_\chi}S.
\]
Thus \(A_d=I-w\otimes w\), which is the inverse of \(\overline g\)
relative to \(g\).  We assume \(p\ge0\), as holds for the profiles used
here.  The formulas involving \(A_\chi\) extend smoothly across
\(a=0\), since
\[
 A_\chi=I-\frac{p+4}{4+pv}\,w\otimes w.
\]
For a symmetric tensor \(T\), let
\[
 p_T=T-(\operatorname{tr}_gT)g,\qquad
 [T_1,T_2]=\langle T_1,T_2\rangle
          -(\operatorname{tr}_gT_1)(\operatorname{tr}_gT_2).
\]
We use the constraint normalization
\[
 16\pi\mu=R_g+6+(\operatorname{tr}_gK)^2-|K|^2,\qquad
 8\pi J=\operatorname{div}_g p_K.
\]

Decompose \(S\) into its \(e^\perp\)-block \(P\), mixed block \(M\),
and axial entry \(q_1=S(e,e)\), and put
\(x_t=\partial_e t\), \(y=(\mathrm dt)_{e^\perp}\).
The quadratic form of Equation~\eqref{eq:ah-T} is
\begin{align}
 \mathcal T={}&\frac12|P^{\mathrm{tf}}|^2
       +|M+(p+1)a y|^2+[4(p+1)-v]|y|^2 \notag\\
 &+4(p+1)d\left(x_t+\frac{\chi a q_1}{4d}\right)^2
       +\frac34\left(F-\frac{\chi q_1}{3}\right)^2
       +\frac{4d(9-d)}{3(4+pv)^2}q_1^2 .
 \label{eq:ah-alg-T}
\end{align}
Although this expression uses an axis where \(a>0\), its sum is
direction independent at \(a=0\), as will also follow from the
coordinate-free expression in the proof below.

\begin{proposition}[Curvature identity]
\label{prop:ah-alg-curvature}
With the preceding definitions, put
\begin{equation}
 V=uA_\chi\bigl(4\nabla Z+K(w,\cdot)^\sharp+A(w)w\bigr).
 \label{eq:ah-alg-V}
\end{equation}
Then the following identity holds pointwise:
\begin{align}
 \frac12e^{4t}(R_{\widehat g}+6)
 ={}&8\pi\bigl(\mu+J(w)\bigr)+3(e^{4t}-1)
       +4\langle A,\mathrm dt\rangle_{A_d}
       +\mathcal T \notag\\
 &+w(F)-F\operatorname{tr}_gK-u^{-1}\operatorname{div}_g V .
 \label{eq:ah-alg-curvature-identity}
\end{align}
Here \(F\) denotes the actual trace
\(\operatorname{tr}_{A_\chi}S\); therefore the identity applies to
any prescribed trace equation by substituting that prescription for
\(F\).
\end{proposition}

\begin{proof}
We first derive the graph identity, keeping the full spatial
dependence of \(l\).  On a product with coordinate \(s\), consider the
stationary Lorentz metric
\[
 \mathbf g=-l^2(\mathrm ds-\mathrm df)^2+\overline g
           =-l^2\mathrm ds^2+2l^2\mathrm ds\,\mathrm df+g .
\]
The \(s\)-slices have induced metric \(g\), shift
\(\beta=l^2\nabla f=Uw\), and lapse \(U=l\sqrt D\).  For the future
normal \(U^{-1}(\partial_s-\beta)\), with the convention
\(k(X,Y)=\mathbf g(\mathbf\nabla_X n,Y)\), stationarity gives
\[
 k=-\frac{1}{2U}\mathcal L_\beta g
   =-H^f-w\otimes B-B\otimes w .
\]
The contracted Gauss equation and the normal Ricci equation for these
stationary slices give
\begin{equation}
 U R_{\mathbf g}
 =U\bigl(R_g+[k,k]\bigr)
   -2\operatorname{div}_g\bigl(\nabla U+(\operatorname{tr}_g k)\beta\bigr).
 \label{eq:ah-alg-stationary}
\end{equation}
For completeness, this sign can be fixed directly by combining
\[
 R_{\mathbf g}=R_g+(\operatorname{tr}k)^2-|k|^2
                   -2\operatorname{Ric}_{\mathbf g}(n,n)
\]
with the normal expansion identity
\[
 \operatorname{Ric}_{\mathbf g}(n,n)
 =-n(\operatorname{tr}k)-|k|^2
       +U^{-1}\Delta_g U .
\]
Indeed \(n(\operatorname{tr}k)=-U^{-1}\beta(\operatorname{tr}k)\)
by stationarity, and
\(\operatorname{div}_g\beta=-U\operatorname{tr}_g k\).
Substitution is exactly
Equation~\eqref{eq:ah-alg-stationary}.

In the coordinate \(s-f\), the same spacetime metric is
\(-l^2\mathrm d(s-f)^2+\overline g\), so
\[
 R_{\mathbf g}=R_{\overline g}-2l^{-1}\overline\Delta l .
\]
The inverse metric and volume element of \(\overline g\) give
\[
 \overline\Delta l
   =D^{-1/2}\operatorname{div}_g(\sqrt D\,A_d\nabla l).
\]
Differentiating \(U\) also gives
\[
 \nabla U-\frac Ul A_d\nabla l=-k(\beta,\cdot)^\sharp.
\]
It follows from Equation~\eqref{eq:ah-alg-stationary} that
\[
 U R_{\overline g}
  =U\bigl(R_g+[k,k]\bigr)+2\operatorname{div}_g(p_k\beta).
\]
Let
\[
 Q=K-k=S+p(w\otimes\mathrm dt+\mathrm dt\otimes w).
\]
Since the symmetric part of \(\nabla\beta\) is \(-Uk\),
\begin{align*}
 \frac2U\operatorname{div}_g(p_K\beta)
  &=2(\operatorname{div}_g p_K)(w)
       +\frac2U\langle p_K,\nabla\beta\rangle\\
  &=16\pi J(w)-2[K,k].
\end{align*}
Using \(p_k=p_K-p_Q\) and
\(R_g=16\pi\mu-6+[K,K]\), we obtain
\begin{equation}
 R_{\overline g}
  =16\pi\bigl(\mu+J(w)\bigr)-6+[Q,Q]
       -2U^{-1}\operatorname{div}_g(Up_Qw).
 \label{eq:ah-alg-graph-curvature}
\end{equation}

The conformal scalar-curvature formula in three dimensions is
\[
 e^{4t}R_{\widehat g}
   =R_{\overline g}-8\overline\Delta t
                         -8|\mathrm dt|_{\overline g}^2 .
\]
Replacing \(U\) by \(u=e^{2t}U\) in the first divergence in
Equation~\eqref{eq:ah-alg-graph-curvature} contributes
\(2p_Q(w,\mathrm dt)\) after division by two.  For the conformal
Laplacian term, the equality
\(\mathrm d\log(e^{2t}l)=(p+2)\mathrm dt+A\) yields
\begin{align*}
 -4\overline\Delta t-4|\mathrm dt|_{\overline g}^2
   ={}&-4u^{-1}\operatorname{div}_g(uA_d\nabla t)\\
     &+4(p+1)|\mathrm dt|_{A_d}^2
       +4\langle A,\mathrm dt\rangle_{A_d}.
\end{align*}
The cosmological terms are \(-3+3e^{4t}\).  Consequently, if
\(V_{\mathrm{pre}}=u(p_Qw+4A_d\nabla t)\), then
\begin{align}
 \frac12e^{4t}(R_{\widehat g}+6)
 ={}&8\pi\bigl(\mu+J(w)\bigr)+3(e^{4t}-1)
       +4\langle A,\mathrm dt\rangle_{A_d}\notag\\
 &+\frac12[Q,Q]+2p_Q(w,\mathrm dt)
       +4(p+1)|\mathrm dt|_{A_d}^2
       -u^{-1}\operatorname{div}_g V_{\mathrm{pre}}.
 \label{eq:ah-alg-before-trace}
\end{align}

We next incorporate the trace equation into the flux.  Direct
differentiation gives
\begin{equation}
 \nabla_iw_j=H^f_{ik}(A_d)_{kj}+dB_iw_j,\qquad
 \mathrm d\log\sqrt D=H^f(w,\cdot)+vB.
 \label{eq:ah-alg-w-derivatives}
\end{equation}
Since \(u=e^{2t}l\sqrt D\), these identities imply
\begin{equation}
 u^{-1}\operatorname{div}_g(uw)
   =F+(1-\chi)q_1-\operatorname{tr}_gK
                   +2(p+1)a x_t .
 \label{eq:ah-alg-uw-divergence}
\end{equation}
One way to see the cancellation of \(A\) is to first obtain
\[
 u^{-1}\operatorname{div}_g(uw)
   =\operatorname{tr}_g H^f+2B(w)+2\mathrm dt(w),
\]
and then use
\(\operatorname{tr}_gS
 =\operatorname{tr}_gK+\operatorname{tr}_gH^f+2A(w)
 =F+(1-\chi)q_1\).
Define \(V=V_{\mathrm{pre}}+uwF\).  The divergence of the added term contributes
\[
 w(F)+F\bigl[F+(1-\chi)q_1-\operatorname{tr}_gK
                         +2(p+1)a x_t\bigr]
\]
to Equation~\eqref{eq:ah-alg-before-trace}.

We verify the resulting quadratic expression explicitly.  Write
\(z=\operatorname{tr}_{e^\perp}P=F-\chi q_1\).  The three blocks
of \(Q\) are \(P\), \(M+pa y\), and \(q_1+2pa x_t\), so
\[
 \frac12[Q,Q]
   =\frac12|P^{\mathrm{tf}}|^2-\frac14z^2
      +|M+pa y|^2-zq_1-2pa x_tz ,
\]
and
\[
 2p_Q(w,\mathrm dt)
    =-2a x_tz+2a\langle M,y\rangle+2pv|y|^2 .
\]
Adding the terms in
Equation~\eqref{eq:ah-alg-before-trace} and
Equation~\eqref{eq:ah-alg-uw-divergence}, and setting aside
\(-F\operatorname{tr}_gK\), gives
\begin{align*}
 &\frac12|P^{\mathrm{tf}}|^2
   +|M+(p+1)a y|^2+[4(p+1)-v]|y|^2\\
 &\quad+4(p+1)d x_t^2+2(p+1)a\chi x_tq_1
       +\frac34F^2-\frac{\chi}{2}Fq_1
       +\left(\chi-\frac{\chi^2}{4}\right)q_1^2 .
\end{align*}
Completing the \(x_t\)-square leaves the coefficient
\[
 \chi-\frac{\chi^2}{4}
      -\frac{(p+1)v\chi^2}{4d}
   =\frac{12d}{(4+pv)^2}
\]
in front of \(q_1^2\).  Completing the \(F\)-square then gives
exactly Equation~\eqref{eq:ah-alg-T}.  Equivalently,
\(\mathcal T\) is the coordinate-free expression
\[
 \frac12[Q,Q]+2p_Q(w,\mathrm dt)
   +4(p+1)|\mathrm dt|_{A_d}^2
   +F\bigl[F+(1-\chi)q_1+2(p+1)a x_t\bigr].
\]
The apparently directional terms in this last expression are smooth:
\((1-\chi)q_1=((p+4)/(4+pv))S(w,w)\) and
\(a x_t=\mathrm dt(w)\).  This proves the asserted extension through
zero gradient.

Finally, Equation~\eqref{eq:ah-alg-w-derivatives} gives
\[
 4\,\mathrm dZ=(4+pv)\mathrm dt+H^f(w,\cdot)+vA.
\]
Using
\((1-\chi)(4+pv)=(p+4)v\), the preceding definition of \(V\)
reduces to
\begin{equation}
 \frac Vu=p_Qw+4A_d\nabla t+wF
       =A_\chi\bigl(4\nabla Z+K(w,\cdot)^\sharp+A(w)w\bigr).
 \label{eq:ah-alg-flux-reduction}
\end{equation}
This proves both the stated flux formula and
Equation~\eqref{eq:ah-alg-curvature-identity}.
\end{proof}

\begin{proposition}[The spatial primitive]
\label{prop:ah-alg-primitive}
Suppose the lapse is the profile defined in
Section~\ref{subsec:ah-profiles}; in particular
\(l_0(t,x)=e^{n(x)t}\) for \(t\le0\) and \(l_0(0,x)=1\).
Define
\begin{gather}
 G(t,x)=\lambda(x)\int_{-\infty}^t L(z,x)\,\mathrm dz,\qquad
 G_c(x)=G(0,x)=\frac{\lambda(x)}{n(x)+2},\notag\\
 V_*=V-4\sqrt D\,A_d\nabla_xG+4\nabla G_c .
 \label{eq:ah-alg-primitive-definitions}
\end{gather}
Assume the scalar equation of Equation~\eqref{eq:ah-system} is written
as
\[
 u^{-1}\operatorname{div}_gV
 =3(e^{4t}-1)+4\langle A,\mathrm dt\rangle_{A_d}+\mathcal P,
\]
where, for the physical system,
\[
 \mathcal P=d_0(\mathcal T+\tau A_Ra\sigma)
            +\rho-m_0 n^{P_0}(\rho_0+v),\qquad \sigma=|\mathrm df|.
\]
Put
\[
 b_1=\frac{G}{\lambda L},\qquad
 s_G=\mathrm d_x\log G,\qquad
 X_1=x_t+\frac{a q_1}{4+pv}.
\]
Then
\begin{equation}
 u^{-1}\operatorname{div}_gV_*=\mathcal P+\mathcal R,
 \label{eq:ah-alg-primitive-divergence}
\end{equation}
where
\begin{align}
 \mathcal R={}&3(e^{4t}-1)-4b_1E_G
           +\frac{4\sqrt d}{\lambda L}\Delta_gG_c\notag\\
 &+4b_1\left\{
       s_G(w)\bigl(F-\operatorname{tr}_{A_d}K+pdaX_1\bigr)
       -pv\langle(s_G)_{e^\perp},y\rangle
       \right\},
 \label{eq:ah-alg-primitive-residual}\\
 E_G={}&\frac{\Delta_xG-\nabla_x^2G(w,w)}{G}
              +v\langle s_G,A\rangle_{A_d}.
 \label{eq:ah-alg-EG}
\end{align}
Here \(\Delta_xG\) and \(\nabla_x^2G\) differentiate \(x\mapsto
G(t,x)\) with \(t\) fixed.  If \(\lambda\) and \(n\) are spatially
constant on a boundary collar, then \(V_*=V\) on that collar.
\end{proposition}

\begin{proof}
For a fixed spatial function \(\phi\), write
\[
 \mathcal L\phi
 =D^{-1/2}\operatorname{div}_g(\sqrt D\,A_d\nabla\phi).
\]
Differentiating \(A_d=I-w\otimes w\) with
Equation~\eqref{eq:ah-alg-w-derivatives} gives
\[
 \mathcal L\phi
    =\operatorname{tr}_{A_d}\nabla^2\phi
        +\mathfrak q(\phi),
\]
where the drift vector is
\begin{equation}
 \mathfrak q
 =vA_d(A+p\,\mathrm dt)^\sharp
  -w\left(F+(d-\chi)q_1-\operatorname{tr}_{A_d}K
                         +2dpa x_t\right).
 \label{eq:ah-alg-primitive-drift}
\end{equation}
To check every contraction, the coefficient of \(\nabla\phi\) before
substituting \(S\) is
\[
 A_d\nabla\log\sqrt D
       -(\operatorname{div}_gw)w-\nabla_ww
 =vA_dB^\sharp
       -w\bigl(\operatorname{tr}_{A_d}H^f+2dB(w)\bigr).
\]
Now
\[
 \operatorname{tr}_{A_d}H^f
  =F+(d-\chi)q_1-\operatorname{tr}_{A_d}K-2dA(w),
\]
which proves Equation~\eqref{eq:ah-alg-primitive-drift}.

In applying this computation to \(\nabla_xG(t(x),x)\), there is one
additional differentiation in \(t\).  The defining integral gives
\[
 G_t=\lambda L,\qquad
 \partial_t\nabla_xG=\nabla_x(\lambda L)=\lambda L A^\sharp.
\]
It follows that
\begin{align*}
 \frac1u\operatorname{div}_g(\sqrt D\,A_d\nabla_xG)
  ={}&b_1\left\{
        \frac{\Delta_xG-\nabla_x^2G(w,w)}G
                 +\langle\mathfrak q,s_G\rangle
        \right\}
      +\langle A,\mathrm dt\rangle_{A_d}.
\end{align*}
Thus the factor \(-4\) in the primitive correction cancels exactly
the prescribed \(4\langle A,\mathrm dt\rangle_{A_d}\).
The \(A\)-part of \(\mathfrak q\) produces \(E_G\).  The remaining
\(\mathrm dt\)-dependent terms are
\[
 4b_1\left\{
   s_G(w)\bigl[(d-\chi)q_1+2dpa x_t\bigr]
                  -pv\langle s_G,\mathrm dt\rangle_{A_d}
       \right\}.
\]
Use
\[
 d-\chi=\frac{pdv}{4+pv},\qquad
 s_G(w)=a(s_G)_e,\qquad
 \langle s_G,\mathrm dt\rangle_{A_d}
       =\langle(s_G)_{e^\perp},y\rangle+d(s_G)_e x_t.
\]
These terms become
\[
 4b_1\left\{s_G(w)pda
                \left(x_t+\frac{a q_1}{4+pv}\right)
                    -pv\langle(s_G)_{e^\perp},y\rangle\right\},
\]
as asserted.  The last correction in \(V_*\) contributes
\(4\Delta_gG_c/u=4\sqrt d\,\Delta_gG_c/(\lambda L)\).
This establishes
Equations~\eqref{eq:ah-alg-primitive-divergence}--\eqref{eq:ah-alg-EG}.
On a collar with spatially constant \(\lambda,n\), both
\(\nabla_xG\) and \(\nabla G_c\) vanish.  The boundary assertion
follows.
\end{proof}

\begin{lemma}[A coercive part of the quadratic form]
\label{lem:ah-alg-coercivity}
For \(p\ge0\) and \(0<d\le1\), the form in
Equation~\eqref{eq:ah-alg-T} satisfies
\begin{equation}
 \mathcal T\ge\frac23F^2+[4(p+1)-v]|y|^2
                    +4(p+1)dX_1^2 .
 \label{eq:ah-alg-coercivity}
\end{equation}
\end{lemma}

\begin{proof}
The first two terms in Equation~\eqref{eq:ah-alg-T} are nonnegative.
The last two terms, keeping the \(X_1\)-square intact, combine to
\[
 \frac34F^2-\frac{2d}{4+pv}Fq_1
                +\frac{12d}{(4+pv)^2}q_1^2.
\]
Completing their square gives the exact identity
\begin{equation}
 \frac34F^2-\frac{2d}{4+pv}Fq_1
       +\frac{12d}{(4+pv)^2}q_1^2
   =\frac{9-d}{12}F^2
      +\frac{12d}{(4+pv)^2}
          \left(q_1-\frac{4+pv}{12}F\right)^2.
 \label{eq:ah-alg-q-minimum}
\end{equation}
Since \((9-d)/12\ge2/3\), the result follows.
\end{proof}

\paragraph{Asymptotic hypotheses for the flux calculation.}
For a tensor \(T\), the notation \(T=O_j(r^{-\gamma})\) means that
its \(b\)-norm and the \(b\)-norms of its covariant derivatives through
order \(j\) are \(O(r^{-\gamma})\).  The following statement concerns
each fixed choice of the profile parameters.  Its constants need not
be uniform as those parameters vary.

\begin{proposition}[Passage to the four metric fluxes]
\label{prop:ah-alg-mass-passage}
Suppose \(g\) has an asymptotically hyperbolic end with background \(b\)
and four finite metric fluxes \(p_a(g)\), \(a=0,1,2,3\), in the
designated chart, as defined in Equation~\eqref{eq:setup-flux}.
Assume that, outside a compact set, \(n\) is a
constant and
\(\lambda=\sqrt{A_0^2+r^2}\), with fixed \(A_0>0\).  Assume the
differentiated end bounds
\begin{equation}
 \begin{gathered}
 g-b=O_2(r^{-\kappa}),\qquad K=O_1(r^{-\kappa}),\\
 t=O_2(r^{-s}),\qquad w=O_2(r^{-\beta}),\\
 \kappa>\frac32,\qquad s>\frac32,\qquad\beta>\frac32,\qquad
 \kappa+\beta>3 .
 \end{gathered}
 \label{eq:ah-alg-end-W}
\end{equation}
Finally, suppose that for some \(\delta>0\),
\begin{equation}
 R_{\widehat g}+6=O(r^{-3-\delta}).
 \label{eq:ah-alg-end-R}
\end{equation}
Then all four metric fluxes of \(\widehat g\) exist and are finite.
For each \(a=0,1,2,3\),
\begin{equation}
 p_a(\widehat g)-p_a(g)
   =-\frac1{2\pi}
       \lim_{r\to\infty}
       \int_{S_r}(V_a\,\mathrm dt-t\,\mathrm dV_a)(\nu_b)
                         \,\mathrm dA_b .
 \label{eq:ah-alg-four-flux}
\end{equation}
Moreover,
\begin{equation}
 V_*=4(\lambda\nabla_gt-t\nabla_g\lambda)+o(r^{-2}),
 \label{eq:ah-alg-Vstar-asymptote}
\end{equation}
and the time-component identity is
\begin{equation}
 p_0(\widehat g)-p_0(g)
   =-\frac1{8\pi}\lim_{r\to\infty}
        \int_{S_r}g(V_*,\nu_g)\,\mathrm dA_g .
 \label{eq:ah-alg-time-flux}
\end{equation}
\end{proposition}

\begin{proof}
We first show that the nonlinear terms have no mass flux.  Since
\(D=(1-|w|_g^2)^{-1}\), the graph correction is exactly
\[
 l^2\mathrm df\otimes\mathrm df
        =D\,w^\flat\otimes w^\flat .
\]
Consequently, with differentiated bounds through order one,
\begin{equation}
 \widehat g-g=4t\,b+\mathcal E,\qquad
 \mathcal E
    =O_1(r^{-2s})+O_1(r^{-s-\kappa})+O_1(r^{-2\beta}).
 \label{eq:ah-alg-metric-error}
\end{equation}
For the four background static potentials,
\[
 |\nabla_b^jV_a|_b=O(r)\quad(j=0,1,2),\qquad
 |S_r|_b=4\pi r^2.
\]
Thus a metric error \(O_1(r^{-j})\) contributes at most
\(O(r^{3-j})\) to its flux.  The assumptions in
Equation~\eqref{eq:ah-alg-end-W} make all three contributions in
Equation~\eqref{eq:ah-alg-metric-error} tend to zero.

Write \(\mathbb U_b(V,e)\) for the one-form
\[
 V(\operatorname{div}_b e-\mathrm d\operatorname{tr}_b e)
       +(\operatorname{tr}_b e)\mathrm dV
       -e(\nabla_bV,\cdot)
\]
appearing in the definition of the metric mass.  In three dimensions,
\begin{equation}
 \mathbb U_b(V,4t\,b)=-8(V\,\mathrm dt-t\,\mathrm dV).
 \label{eq:ah-alg-conformal-flux}
\end{equation}
Indeed \(\operatorname{div}_b(4tb)=4\mathrm dt\),
\(\operatorname{tr}_b(4tb)=12t\), and
\((4tb)(\nabla_bV,\cdot)=4t\,\mathrm dV\).  Combining
Equation~\eqref{eq:ah-alg-conformal-flux} with
Equation~\eqref{eq:ah-alg-metric-error} proves the asserted
comparison of the finite-radius flux integrals.

We establish existence of the limiting fluxes before taking their
difference.  Put \(\widehat e=\widehat g-b\) and
\(\nu=\min(\kappa,s,2\beta)>3/2\).  Then
\(\widehat e=O_2(r^{-\nu})\).  The scalar-curvature linearization at
the curvature \(-1\) background is
\[
 D R_b(e)=\operatorname{div}_b\operatorname{div}_b e
           -\Delta_b\operatorname{tr}_b e+2\operatorname{tr}_b e.
\]
Since \(\nabla_b^2V_a=V_a b\) and \(\Delta_bV_a=3V_a\),
differentiating the defining one-form gives exactly
\begin{equation}
 \operatorname{div}_b\mathbb U_b(V_a,\widehat e)
                   =V_a D R_b(\widehat e).
 \label{eq:ah-alg-linearized-divergence}
\end{equation}
The Taylor remainder in scalar curvature has the estimate
\[
 \bigl|D R_b(\widehat e)-(R_{\widehat g}+6)\bigr|
 \le C\bigl(
       |\widehat e|\,|\nabla_b^2\widehat e|
       +|\nabla_b\widehat e|^2+|\widehat e|^2
       \bigr)
 =O(r^{-2\nu}).
\]
The constant is uniform on a sufficiently distant tail, where
\(\widehat g\) and \(b\) are uniformly equivalent.  As
\[
 \mathrm dV_b=\frac{r^2}{\sqrt{1+r^2}}\,\mathrm dr\,\mathrm d\omega,
\]
the scalar-curvature contribution and the Taylor remainder in the
right side of Equation~\eqref{eq:ah-alg-linearized-divergence} are
absolutely integrable there:
\[
 \int^\infty r^{-1-\delta}\,\mathrm dr<\infty,\qquad
 \int^\infty r^{2-2\nu}\,\mathrm dr<\infty .
\]
The divergence theorem on annuli therefore makes each family of
flux integrals a Cauchy family.  All four limits exist and are finite.
Taking the difference from the given fluxes of \(g\) now proves
Equation~\eqref{eq:ah-alg-four-flux}.

It remains to identify the corrected vector field with the time
potential.  On the ultimate tail, set
\[
 H_n(t)=\int_{-\infty}^t L(z)\,\mathrm dz.
\]
The constant \(n\) is fixed here, and
\[
 H_n(0)=\frac1{n+2},\qquad H_n'(0)=1,\qquad
 \nabla_xG=H_n(t)\nabla_g\lambda .
\]
Taylor expansion at zero gives
\[
 \nabla_xG
   =\left(\frac1{n+2}+t+O(t^2)\right)\nabla_g\lambda,\qquad
 \nabla_gG_c=\frac{\nabla_g\lambda}{n+2}.
\]
On the same tail \(A=\mathrm d\log\lambda\), so
\[
 \sqrt D A_d=I+O(r^{-2\beta}),\qquad
 \frac u\lambda=1+O(r^{-s})+O(r^{-2\beta}),\qquad
 \nabla Z=\nabla t+O(r^{-2\beta}).
\]
Substituting these bounds into
Equation~\eqref{eq:ah-alg-V} and
Equation~\eqref{eq:ah-alg-primitive-definitions} yields the explicit
error ledger
\begin{equation}
 V_*-4(\lambda\nabla_gt-t\nabla_g\lambda)
   =O(r^{1-2s})+O(r^{1-2\beta})
                         +O(r^{1-\kappa-\beta}).
 \label{eq:ah-alg-vector-error}
\end{equation}
All three terms are \(o(r^{-2})\) by
Equation~\eqref{eq:ah-alg-end-W}.  In particular their integrated
normal fluxes tend to zero.

Finally,
\[
 \lambda-V_0=O_1(r^{-1}),\qquad V_0=\sqrt{1+r^2}.
\]
Replacing \(V_0\) by \(\lambda\) in the conformal flux integral costs
\(O(r^{1-s})=o(1)\).  Replacing \(\nu_b,\mathrm dA_b\), and the
background gradient by their \(g\)-counterparts costs
\(O(r^{3-s-\kappa})=o(1)\), since \(s+\kappa>3\).
Equation~\eqref{eq:ah-alg-time-flux} follows from
Equations~\eqref{eq:ah-alg-four-flux} and
\eqref{eq:ah-alg-Vstar-asymptote}.  The proof treats the four static
potentials separately; it does not assume a causal character for the
mass covector of \(\widehat g\).
\end{proof}

\begin{lemma}[Scalar-curvature decay from the end equation]
\label{lem:ah-alg-scalar-decay}
Suppose the differentiated bounds of
Equation~\eqref{eq:ah-alg-end-W} hold,
\(R_g+6=O(r^{-3-\delta})\), and
\[
 2s>3+\delta,\qquad 2\beta>3+\delta.
\]
Suppose also that the scalar end equation has the form
\begin{equation}
 \Delta_gZ+C(Z)|\mathrm dZ|_g^2
      =\frac34(e^{4Z}-1)+O(r^{-3-\delta}),
 \label{eq:ah-alg-end-equation}
\end{equation}
where \(C\) is smooth and bounded near zero for the fixed profiles.
Then Equation~\eqref{eq:ah-alg-end-R} holds.  In the deformation
equations one takes
\[
 C(Z)=p(Z)+2-d_0(Z)\bigl(p(Z)+1\bigr)
\]
on the tail with spatially constant \(n\).
\end{lemma}

\begin{proof}
Since \(D=(1-|w|^2)^{-1}\), we have
\(Z-t=O_2(r^{-2\beta})\).  The assumed end equation therefore gives
\[
 \Delta_gt+C(t)|\mathrm dt|_g^2
      =\frac34(e^{4t}-1)+O(r^{-3-\delta}).
\]
First apply the conformal scalar-curvature formula to \(e^{4t}g\):
\begin{align*}
 e^{4t}\bigl(R_{e^{4t}g}+6\bigr)
   &=R_g+6-8\Delta_gt-8|\mathrm dt|_g^2+6(e^{4t}-1)\\
   &=R_g+6+8\bigl(C(t)-1\bigr)|\mathrm dt|_g^2
                                      +O(r^{-3-\delta}).
\end{align*}
The gradient term is \(O(r^{-2s})=O(r^{-3-\delta})\).
The difference
\(\widehat g-e^{4t}g=e^{4t}D\,w^\flat\otimes w^\flat\)
is \(O_2(r^{-2\beta})\).  The local scalar-curvature formula,
on the uniformly equivalent end metrics with bounded curvature,
then gives
\(R_{\widehat g}-R_{e^{4t}g}=O(r^{-2\beta})\).
This proves the required decay.
\end{proof}

\begin{remark}
The scalar-curvature hypothesis in the mass proposition is essential
to the convergence argument.  Bounds \(t=O_2(r^{-s})\) for every
\(s<3\), by themselves, would allow \(t=r^{-3}\log r\), whose
conformal time-flux integral diverges logarithmically.  In the
construction, Lemma~\ref{lem:ah-alg-scalar-decay} supplies the missing
integrability from the scalar equation.  The estimates of the
primitive residual are then used separately to determine the sign of
the limiting flux in Equation~\eqref{eq:ah-alg-time-flux}.
\end{remark}
